\documentclass[a4paper]{article}

\usepackage[utf8]{inputenc}
\usepackage[spanish]{babel}
\usepackage{amsmath,amssymb,amsfonts}
\usepackage{graphicx}
\usepackage{enumitem}

% Definir los entornos necesarios para exams
\newenvironment{question}{}{}
\newenvironment{solution}{}{}
\newenvironment{answerlist}{\begin{enumerate}}{\end{enumerate}}

\usepackage{Sweave}
\begin{document}
\Sconcordance{concordance:boxplots.tex:boxplots.Rnw:1 13 1 1 0 2 1 1 126 7 1 1 4 2 2 %
5 0 1 2 3 1 1 4 5 0 1 2 7 1}



\begin{question}
En la siguiente figura se representan las distribuciones de una variable
dada por dos muestras (A y B) mediante diagramas de caja paralelos.
¿Cuáles de las siguientes afirmaciones son correctas? \emph{(Comentario: Las
afirmaciones son correctas o claramente incorrectas.)}

\setkeys{Gin}{width=0.7\textwidth}
\includegraphics{boxplots-002}

\begin{answerlist}
  \item La ubicación de ambas distribuciones es aproximadamente la misma.
  \item Ambas distribuciones no contienen valores atípicos.
  \item La dispersión en la muestra A es claramente mayor que en B.
  \item La asimetría de ambas muestras es similar.
  \item La distribución A es sesgada a la izquierda.
\end{answerlist}
\end{question}

\begin{solution}
\begin{answerlist}
  \item Falso. La distribución B tiene en promedio valores más altos que la distribución A.
  \item Verdadero. Ambas distribuciones no tienen observaciones que se desvíen más de 1.5 veces el rango intercuartílico desde la caja.
  \item Verdadero. El rango intercuartílico en la muestra A es  claramente mayor que en B.
  \item Verdadero. La asimetría de ambas distribuciones es similar, ambas son aproximadamente simétricas.
  \item Falso. La distribución  A  es  aproximadamente simétrica.
\end{answerlist}
\end{solution}

%% META-INFORMATION
%% \extype{mchoice}
%% \exsolution{01110}
%% \exname{Diagramas de caja paralelos}
\end{document}
